% ═══════════════════════════════════════════════════════════════════
\section{Why the Search Does Not Localise the Binding Site}
\label{sec:why}
% ═══════════════════════════════════════════════════════════════════

Section~\ref{sec:results} shows that the pipeline does not find the binding site. A
negative benchmark is only useful if we can say which part of the system is responsible,
and there are three candidates: the classical stages that build the representation, the
encoding that turns a window of that representation into a quantum state, or the search
that queries the encoding. This section argues that the encoding is responsible. We first
rule out the explanations our measurements do not support, then describe what the
descriptor actually contains, and finally say what the result does and does not tell us
about the quantum part.

\subsection{Ruling out the alternatives}
\label{sec:why-not}

The first explanation to consider is that we rank the candidates badly, and at first
sight the benchmark supports it, since the candidate we rank first is on average worse
than the average of the candidates we return (47.2nd percentile against 32.2nd,
Table~\ref{tab:percentile}). Table~\ref{tab:ksweep} shows that this reading is wrong. At
$k=3$ the whole set of windows matching the ligand's bitstring sits at the 48.2nd
percentile, which is chance. A ranking function only reorders a set, so it cannot add
information the set does not already contain, and any criterion applied to that set has
the expected accuracy of a random pick. The poor showing of \texttt{interactivity\_score}
is therefore not evidence that this particular score is bad: when the underlying set is
at chance, an ordering is as likely to come out anti-correlated with the truth as
correlated with it. Ranking is not the bottleneck because there is nothing to rank.

The matching criterion is not the bottleneck either, although it does cause a different
problem that the same experiments let us separate out. Replacing exact bitstring equality
with a Euclidean distance in second encoding amplitude space raises coverage from 6 to 27
of 29 pockets at $k=5$, which is a real improvement and the one countermeasure of the two
that worked, but accuracy stays at the 45.8th percentile against 50\% for chance
(Section~\ref{sec:res-sweep}). Window size behaves the same way: the improvement we see
at $k=5$ under exact matching (21.5\% against a 38.2\% baseline) comes with the collapse
in coverage that pushed us to the amplitude variant in the first place, and it does not
survive once coverage is restored. Coverage and accuracy therefore have different causes,
and only coverage depends on how we match.

The result is also not caused by a few bad targets, nor by the defects we list in
Section~\ref{sec:limitations}. The Spearman correlation between amplitude space
similarity to the ligand and true spatial distance exceeds $|\rho| > 0.3$ in only 2 of 29
pockets, and its sign is correct in 16 of 29, barely more than half, as we would expect
from noise (Section~\ref{sec:res-spearman}). The outcome is uniform across the benchmark
rather than an average of successes and failures. The implementation defects are a more
serious objection, since the largest of them puts spurious \texttt{PosIonizable} sites on
13.5\% of the features we extract, and they are worth fixing regardless. They change the
values carried in the $hb$ channel, though, and an error of that size degrades a
correlation rather than taking it from something useful down to $\rho \approx 0$. It
also does not touch the window geometry figures of Section~\ref{sec:res-geometry}, which
are computed from coordinates only.

That leaves the explanation we found most convincing when we started, which is that
flattening the pocket from 3D to a 1D chain destroys the spatial arrangement of the
sites, so that a window of consecutive positions is an arbitrary collection with no
counterpart on the protein surface. Our own measurements do not support it. At $k=5$ the
median window span is 7.63~\AA{} and the median radius of gyration is 3.09~\AA{}, and only
25.8\% of windows are wider than 10~\AA{}. Consecutive sites do form compact local
patches, roughly the size of a small ligand, which means chain assembly and topological
ordering largely do what we designed them to do. A good deal of the geometric information
is therefore still there, implicitly, in the windows the search looks at, and the
descriptor fails on windows that are spatially sensible. So the explanation has to be in
what a window is described by, not in how a window is put together.

\subsection{Composition alone does not identify a binding site}
\label{sec:why-mechanism}

What is left, after the eliminations above, is that a window descriptor records the local
pharmacophoric composition of a patch of surface: which interaction families are present
there and how strong they are. The problem is that this composition is close to the same
everywhere, and our own extraction statistics say so. \texttt{HBondDonor} and
\texttt{HBondAcceptor} appear in 94\% to 100\% of residues, because every amino acid
contributes a backbone nitrogen and a backbone oxygen; \texttt{Hydrophobe} appears in
79\% to 84\%; and the average number of features per residue holds between 4.8 and 5.2 across
structures ranging from 105 to 1239 residues (Table~\ref{tab:featval}). We reported those
numbers in Section~\ref{sec:feature-extraction} as evidence that extraction is stable and does not
depend on protein size, which they are. They also mean that the same signature turns up
all over the surface, so a descriptor built from composition gives similar values to
patches that have nothing to do with each other. The search is not looking in the wrong
way so much as looking for something that is not rare.

The first encoding makes this worse, and it does so by construction rather than by
accident. Reducing each site to two bits gives a window one of $4^k$ possible values, 64
at the default $k=3$, while a single pocket produces somewhere between roughly 100 and
600 windows. Many windows must therefore share a code, and nothing in how we build the
code makes the windows that share one spatially related. The same arithmetic explains the
coverage numbers in Table~\ref{tab:ksweep}: raising $k$ grows the code space faster than
it grows the number of windows, so exact matches become rare without becoming selective,
which is why coverage drops to 7 of 29 pockets at $k=5$ while accuracy barely moves.

There is a second loss, and this one is about which distinctions the encoding can express
at all. The reduction to two channels in Section~\ref{sec:quantum-encoding} sends
\texttt{HBondDonor}, \texttt{HBondAcceptor}, \texttt{PosIonizable} and
\texttt{NegIonizable} all to the same $hb$ channel, so a donor and an acceptor, and
equally a positive and a negative group, look identical in the descriptor. Binding works
by complementary groups pairing up: a donor on the ligand is satisfied by an acceptor on
the protein, not by another donor. An encoding that cannot tell those two apart cannot
represent the pairing that decides whether a ligand binds, so asking the oracle for a
window whose code is equal to the ligand's is asking the wrong question, independently of
everything above.

\subsection{What this says about the quantum part, and what would have to change}
\label{sec:why-forward}

It matters, for how this result gets cited, that none of the above is an argument about
Grover search. Our oracle is
$\hat{O} = I - 2\ket{x_{\text{lig}}}\!\bra{x_{\text{lig}}}$, which tests equality against
a fixed bitstring, so the set of windows that satisfy it is fixed by the encoding before
any circuit is built. Amplitude amplification changes which element of that set we
measure, never which windows are in it. The quantum stage speeds up a lookup; it does not
score candidates, and it can neither improve nor damage spatial accuracy. Our negative
result is therefore not evidence that a quantum approach to docking site identification
does not work. The quadratic speed-up is real, and we applied it to a dictionary whose
keys do not tell binding sites apart from the rest of the surface. Strictly speaking the
search stage is still untested, because an algorithm cannot be evaluated on a problem
where every candidate is equally good.

The eliminations in Section~\ref{sec:why-not} also narrow down what a redesign would have
to provide, and this is where a negative result becomes useful. Since windows are already
compact, adding the distances between the sites of a window, which is what classical
pharmacophore fingerprints encode as combinations of types and distances, would mostly
write down information that is already implicitly there, and we would not expect it on
its own to change the outcome. What separates a binding pocket from any other patch of
surface is not how a handful of sites are arranged but how enclosed they are. A pocket is
a cavity, and what makes it able to host a ligand is the empty space inside it, which no
list of sites and intensities describes at any window size.

Our report already names this quantity, in a different context. Section~\ref{sec:surface-filter}
notes that the SASA filter measures solvent accessibility globally and cannot tell an
atom facing the inside of a pocket from one exposed on the outer surface of the protein,
and that tools such as fpocket and CASTp exist precisely to compute which atoms line a
detected cavity. We wrote that as a limitation of one filter. Read together with
Section~\ref{sec:results}, it is better understood as naming what the whole
representation leaves out: our pipeline describes where the atoms are and never where a
ligand would fit.

An encoding able to support this search would therefore need at least three things: some
measure of how buried or enclosed a site or a candidate region is; separate channels for
the hydrogen bonding and ionic families, so that complementarity can be expressed at all;
and candidate regions built from spatial neighbourhoods instead of windows of the
sequence, so that residues on opposite walls of the same pocket can end up in one
descriptor. All three are compatible with the two step encoding and with the same Grover
formulation downstream. That is the sense in which our result is a specification rather
than a dead end: it does not show that the scheme of the reference paper is wrong, it
identifies what has to be put into the interaction space before that scheme has anything
to search.


% ═══════════════════════════════════════════════════════════════════
\section{Known Limitations}
\label{sec:limitations}
% ═══════════════════════════════════════════════════════════════════

This section collects the problems we know about in the current implementation. We group
them into three parts: problems that come from reading real PDB files, defects we found
by running our own validation scripts over the whole benchmark, and limitations of the
model and of the libraries we used. The last subsection separates all of these from the
limitation that actually explains the benchmark result of Section~\ref{sec:why}.

\subsection*{Template matching ambiguities and missing atoms}

Step~1 assigns bond orders by matching each residue against a template. This produces two
kinds of diagnostic message across the four test structures.

The first is a warning from RDKit, \texttt{More than one matching pattern found}. It
appears when a residue contains two atoms that are interchangeable as far as the template
is concerned, so there is more than one valid way to match it and RDKit picks one of them
arbitrarily. It happens on glutamate, aspartate and arginine. The choice does not matter
for us: the two matches differ only in which of the two interchangeable atoms is used,
and the pharmacophoric sites we extract afterwards come out identical either way. The
number of warnings grows with the size of the protein, from 43 on 1a1c (221 residues) to
249 on 1a0t (1239 residues).

The second kind of message concerns residues that matched no template at all and were
skipped. Table~\ref{tab:step1-validation} reports how many.

\begin{table}[h]
\centering
\caption{Step~1 validation results across four protein structures. Skipped residues are
those where bond order assignment failed because the file does not contain all the atoms
the template expects.}
\label{tab:step1-validation}
\begin{tabular}{lrrrr}
\toprule
\textbf{Protein} & \textbf{Total residues} & \textbf{Processed} &
\textbf{Skipped} & \textbf{Skip rate} \\
\midrule
1a08 & 105  & 99   & 6  & 5.7\% \\
1a0t & 1239 & 1239 & 0  & 0.0\% \\
1a1c & 209  & 197  & 12 & 5.7\% \\
1a46 & 280  & 272  & 8  & 2.9\% \\
\bottomrule
\end{tabular}
\end{table}

The skip rate goes from 0\% to 5.7\%. Looking at which residues get skipped, they are the
ones whose record in the file is missing some of the atoms the template requires, which
is common in experimentally determined structures published without further processing.
1a0t is the largest of the four structures and skips nothing at all, so the failures
depend on how complete the input file is and not on the size of the protein or on our
code.

The usual fix is to run the input through a structure completion tool such as
\texttt{PDBFixer}, which fills in the missing atoms before our pipeline starts. We did
not do this, and of all the fixes listed in this section it is the easiest to add.

\subsection*{Defects found by validating the full benchmark}

The table above covers four structures. Running our five validation scripts
(\texttt{scripts/validate\_step1.py} to \texttt{validate\_step5.py}) over all 31 benchmark
pockets instead brings out four defects that appear consistently, and that a sample of
four structures was too small to reveal. Table~\ref{tab:systematic-defects} lists them and
the rest of this subsection explains each one.

\begin{table}[h]
\centering
\caption{Defects found by running the validation scripts over all 31 benchmark pockets,
and their state after the fixes described below. \emph{When found} and \emph{now} are the
same measurement repeated on the same 31 pockets.}
\label{tab:systematic-defects}
\begin{tabular}{lllr}
\toprule
\textbf{Defect} & \textbf{Cause in the code} & \textbf{When found} & \textbf{Now} \\
\midrule
Sites marked \texttt{PosIonizable} by mistake & residues built one at a time & 1034/7675 features (13.5\%) & 1/5247 (0.02\%) \\
No residue ever carries a charge              & templates describe uncharged residues & ARG 66/66, ASP 57/57, & every one ionised \\
                                              &                                       & GLU 55/55, LYS 23/23  & \\
Deduplication does not converge               & greedy grouping, no second pass       & 24/31 pockets & 0/31 observed\textsuperscript{\dag} \\
Insertion codes dropped                       & residue key ignores column 27         & 4 records, 1YGC and 1OYT & open \\
\bottomrule
\end{tabular}

\vspace{2pt}
{\footnotesize \textsuperscript{\dag}No longer observed on the benchmark; the algorithm
is still single-pass and a counterexample exists (see below).}
\end{table}

The first two defects have the same cause. We build each residue on its own, rather than
as a link in a chain, so every residue is left with an open end at the point where it
would normally be attached to the previous one. RDKit reads that open end as a group that
can take a positive charge and creates a \texttt{PosIonizable} site for it. In a real
protein chain the site is not there, and these false sites made up 89.4\% of all the
\texttt{PosIonizable} sites we produced. The same simplification caused the second
defect: the residue templates in \texttt{pdb\_reader.py} described uncharged versions of
the residues, so nothing we output ever carried an electric charge, even though several
residues are normally charged in a real protein.

Both are now fixed. \texttt{extract\_features} discards any \texttt{PosIonizable} whose
atoms are a backbone amide nitrogen, and the templates carry the ionised side chains
(\texttt{[NH3+]} on LYS, \texttt{[NH2+]} on ARG, \texttt{[O-]} on ASP and GLU; HIS is
left neutral, as it mostly is at physiological pH). Repeating the same count over the
same 31 pockets leaves 1 spurious \texttt{PosIonizable} out of 5247 features, and the
122 that remain sit on the side chains of ARG, LYS and HIS, where they belong.

The third defect is in \texttt{site\_dedup\_centroid.py}. The algorithm merges nearby
sites of the same type by grouping them greedily and replacing each group with its
centroid. The problem is that the centroid keeps moving as members are added to the
group, and we never went back to check the earlier members against the final position, so
Test~1 of \texttt{validate\_step4.py} failed on 24 of the 31 pockets. The sites involved
are nearly always a pair of atom groups in valine and leucine that sit about 1.38~Å
apart, just below our 1.5~Å merging threshold. Comparing each site against the group's
running centroid, rather than against its first member, removes the failure on this
benchmark: Test~1 now passes on 31 of the 31 pockets. It does not make the algorithm
correct in general. The grouping is still a single greedy pass with no second sweep over
the final centroids, and three collinear sites of the same type at $0$, $1.6$ and
$1.4$~Å, presented in that order, still produce two centroids $0.9$~Å apart --- below the
threshold the step is supposed to enforce. The result also depends on the order in which
the sites arrive: the same three sites given as $0$, $1.4$, $1.6$ merge into one. What we
can claim is that the benchmark no longer exhibits the defect, not that it cannot occur.
Agglomerative clustering, or repeating the assignment until it stops changing, would
settle it.

The fourth defect is a parsing bug. The PDB format reserves column 27 of an \texttt{ATOM}
record for an insertion code, which lets two different residues share the same sequence
number, numbered for instance 170 and 170A. \texttt{load\_residues\_from\_pdb} reads only
columns 23 to 26 and uses $(\texttt{chain\_id}, \texttt{res\_seq}, \texttt{res\_name})$ as
its dictionary key, so in a file that uses insertion codes two different residues that
happen to have the same name end up in one record holding the atoms of both. We found
this in two benchmark files, on four records in total: in 1YGC, H60\_LYS holds 18 atoms
where a lysine has 9, H170\_GLN holds 18 where a glutamine has 9, and H170\_SER holds 12
where a serine has 6; in 1OYT, H60\_PRO holds 14 where a proline has 7. None of them is
caught downstream --- the merged molecule passes bond-order assignment without a warning
and goes on to produce pharmacophoric sites.

This defect is also the cause of the only other failure the validation reports. Test~3 of
\texttt{validate\_step4.py} checks that deduplication conserves total intensity, and it
fails on exactly one pocket of the 31: 1YGC, on exactly the two records above
(H60\_LYS loses 0.4323, H170\_GLN loses 0.1441). The mechanism is the merge. A doubled
residue yields 15 distinct sites where a real one yields 8 or 9, which is more than the
12 that \texttt{select\_dedup\_centroid} keeps, so three sites are dropped and their
intensity with them. No correctly parsed residue in the benchmark comes near the cap.
Fixing the key therefore removes both failures; the site cap is not independently at
fault. The fix is to add the insertion code to the key and to
\texttt{ResidueRecord}, and it is not implemented.

Two remarks on how to read these results. First, the same run confirms that feature
extraction passes on 31/31 pockets with no residue missing a pharmacophore type it should
have and no residue producing zero features, that no site gets a negative intensity, and
that the BFS ordering is deterministic on 31/31. The defects above are localised rather
than spread over the whole pipeline. Second, one reported failure is not a defect: Test~3
of \texttt{validate\_step5.py} fails on 31/31 pockets, but the test is written
incorrectly. It compares site sequences after the SASA filter has run, and that filter
removes a different number of sites from each copy of a residue type, so the sequences
cannot be expected to match. Measured before the filter, the inconsistencies drop from
141 to 47 and disappear for the residues whose shape is fixed. The ones left are on
residues that can bend, and follow from deduplication depending on the shape the residue
happens to take, which is the intended behaviour.

\subsection*{Limitations of the model and of the libraries}

\begin{itemize}[leftmargin=1.5em]
    \item \textbf{Intensity values for charged residues.} The intensities we assign to
    hydrogen bonding sites come from the Abraham scales, which were measured for
    uncharged molecules in solution. Several side chains (lysine, arginine, aspartate,
    glutamate) normally carry a charge, and for those we substitute the value of the
    closest uncharged molecule available, which we expect underestimates how strongly
    they really interact. The effect shows up in our own output: arginine and lysine,
    usually described as strong hydrogen bond donors, receive lower values (around 0.22
    to 0.28) than serine and the acidic residues (around 0.70), because the scale rates
    the bond on its own and ignores the charge. The intensities we assign to these
    residues should therefore be read with caution. Doing better would need scales
    measured for the charged forms, which are not in the standard tables, or a different
    way of computing intensity altogether.
    \item \textbf{Intensity coverage for ligands.} Our intensity table is indexed by
    residue name and atom name, so it covers protein residues and not ligand atoms. The
    fallback recognises only four groups by name (carboxyl, amine, phenol, thiol), and
    anything else keeps the neutral default value.
    \item \textbf{Qubit count does not scale.} Qiskit builds our oracle and diffusion
    operators from dense $2^{n} \times 2^{n}$ matrices, and its general purpose unitary
    synthesis does not scale past roughly 10 qubits. This caps
    \texttt{--ligand-max-sites} at small values (default 3, so 6 qubits) if we want the
    pipeline to finish in seconds rather than minutes, and it is why the sweep over
    window sizes in Section~\ref{sec:res-sweep} is computed classically instead of by
    running the circuit for every configuration.
    \item \textbf{SWAP test not implemented as a circuit.} We did not implement the SWAP
    test proposed by the reference paper. We did implement and evaluate its classical
    equivalent, a Euclidean distance between the amplitude vectors of the second
    encoding, reported in Section~\ref{sec:res-sweep}: it raises coverage from 6 to 27 of
    29 pockets and does not improve accuracy. Since the quantum stage does not change
    which windows get selected (Section~\ref{sec:why-forward}), we would expect a circuit
    implementation to give the same accuracy.
    \item \textbf{Modules not wired into the pipeline.} \texttt{lattice\_fitting.py},
    \texttt{snapping.py} and \texttt{labeling.py} implement a 3D to 2D PCA projection,
    integer lattice snapping and one-hot pharmacophore labelling. They are written and
    tested but never called from \texttt{run\_pipeline.py}.
\end{itemize}

\subsection*{Why none of this explains the benchmark result}

Everything above is a defect or a shortcut in the implementation, and all of it can be
fixed without changing the design. None of it explains the result of
Section~\ref{sec:why}. These problems change the values attached to individual sites,
while what fails in the benchmark is the relationship between how similar two descriptors
are and how far apart they are in space, and that relationship is absent rather than
noisy. The window geometry figures in Section~\ref{sec:res-geometry} do not use
intensities at all, only coordinates, so none of the defects listed here affects them.
Fixing these problems is worth doing, and it has to happen before any redesign can be
evaluated, but on its own it would not change the outcome.

The limitation that does explain the outcome is in the design and not in the code. We
describe a candidate docking site by which interaction types appear in a fixed length
window of sites, and that description leaves out two things: how enclosed the space is
that a ligand would have to occupy, and whether the ligand's groups and the protein's
groups fit together. No choice of window size, ranking criterion, threshold or matching
rule compensates for this, because neither quantity is represented at any setting of
those parameters. A narrower consequence of the same choice is that a candidate region is
a window of the sequence, so two sites that are close in space but far apart along the
chain, as the sites on opposite walls of a pocket usually are, can never appear in the
same descriptor.


% ═══════════════════════════════════════════════════════════════════
\section{Conclusion}
\label{sec:conclusion}
% ═══════════════════════════════════════════════════════════════════

Q-MODE implements a complete path from a PDB binding pocket to a Grover search over
candidate docking windows executed on a quantum circuit, with site intensities taken from
measured physical quantities, Crippen contributions for the hydrophobic channel and
Abraham scales for the hydrogen bonding channel, rather than from placeholder constants. We validated the classical stages across all 31 benchmark pockets: feature
extraction passes on 31/31 with no residue missing an expected interaction type and no
residue producing zero features, no site is given a negative intensity, and the BFS
topological ordering is deterministic on 31/31. The encoding and search stages reproduce
the first encoding, second encoding, oracle and diffusion operator of the reference
paper.

The end to end evaluation is negative, and that is the main result of this report. Across
29 benchmark pockets the candidate windows the search returns are no closer to the
crystallographic ligand than windows picked at random from the same pocket: the best
candidate sits at the 32.2nd percentile against a 30.8\% random baseline, and the top
ranked candidate at the 47.2nd against 50\% for chance. The same result in absolute units,
which is how it is most easily read: the closest window that exists in the representation
is a median 3.76~\AA{} from the true ligand centroid, drawing a window at random gives
7.88~\AA{}, and the search returns 7.75~\AA{}. Of the 4.12~\AA{} that separated chance
from the best answer available to it, it recovers 0.13~\AA{} --- which is to say
essentially none. That answer was available: in
25 of the 29 pockets a window within 5~\AA{} of the ligand centroid is physically present
in the chain, so the failure is not an artefact of pocket geometry or of an unreachable
success criterion. Neither of the two countermeasures we tried changes this. A wider window improves the percentile of the matching set but
drops coverage to 7 of 29 pockets, and replacing exact bitstring equality with a distance
in amplitude space brings coverage back to 27 of 29, which is a genuine improvement and
the part worth keeping, while leaving accuracy at chance.

The cause is in the representation, upstream of the quantum stage. A window descriptor
records the pharmacophoric composition of a patch of surface, and that composition is
close to uniform: the backbone alone puts a donor and an acceptor on 94\% to 100\% of
residues, so the same signature recurs everywhere. The reduction to two channels also
merges donors with acceptors and positive with negative groups, which removes the
complementarity that decides whether a ligand binds. One thing our measurements do not
support is the explanation we expected to find. The failure cannot be blamed on losing
the spatial arrangement of the sites when we flatten the pocket into a chain, because the
windows turn out to be compact patches, with a median span of 7.63~\AA{} and a radius of
gyration of 3.09~\AA{}, so their local geometry is largely preserved and the descriptor
fails on windows that are spatially sensible. What the representation never encodes is
how enclosed a region is, which is the property that separates a pocket from any other
piece of surface, and the one that Section~\ref{sec:surface-filter} already noted a globally
computed SASA cannot see.

Because the oracle tests equality and the set of matching windows is fixed classically
before any circuit runs, this result constrains the interaction space encoding and leaves
the quantum search itself untested: an algorithm cannot be judged on a problem where
every candidate is equally good. Read that way the outcome is a specification rather than
a dead end. It says that an encoding meant to support docking site search has to carry
some measure of burial or enclosure, has to keep the hydrogen bonding and ionic families
separate so that complementarity can be expressed, and has to build candidate regions
from spatial neighbourhoods rather than from windows of the sequence. All of these fit
inside the two step encoding scheme and the Grover formulation of the reference paper.
The defects we list in Section~\ref{sec:limitations} are, for the most part, behind us:
the spurious \texttt{PosIonizable} sites and the missing charge state are fixed, and the
deduplication no longer misbehaves anywhere on the benchmark, though the algorithm that
produces it is still the single greedy pass and can be made to fail by construction. The
benchmark reported here was re-run after those fixes. The dropped insertion codes remain
open, on four records across two files. That the
numbers did not move is itself part of the result: the defects were real, and correcting
them changed nothing about the outcome.
